\section{Collective content, primitive norm profiles, and exact complement cancellation}
\label{sec-cc-collective-profile}

Throughout this section let
\[
 \mathcal O=\mathbb Z[\zeta],\qquad \zeta^2-\zeta+1=0,\qquad
 N(a+b\zeta)=a^2+ab+b^2.
\]
For a nonzero \(Z=A+C\zeta\), put
\[
 \operatorname{cont}(Z)=\gcd(|A|,|C|)>0,
 \qquad Z_{\rm prim}=Z/\operatorname{cont}(Z).
\]
Fix an integer \(B\ge2\), and define the actual input block
\[
 K_B=\{B,\ldots,2B-1\},\qquad w_k=3k+\zeta,
 \qquad Q_k=N(w_k)=9k^2+3k+1.
\]
The results below concern these input norms. They do not identify their
prime factors with the boundary primes of \(P(w_k^n)\), nor assume that
the whole block of ratios \(w_k/\bar w_k\) is multiplicatively independent.

\subsection{The exact oriented content and its interval average}

Each \(Q_k\) is odd, is one modulo three, and has only split prime
factors \(p\equiv1\pmod3\). For each such prime choose the two conjugate
prime ideals \((\pi_p),(\bar\pi_p)\), and write
\[
 d_{k,p}=v_{\pi_p}(w_k)-v_{\bar\pi_p}(w_k).
\]
The two valuations cannot both be positive, since the coordinates
\((3k,1)\) are primitive. Hence \(|d_{k,p}|=v_p(Q_k)\).

\begin{lemma}[CC1: exact signed incidence and content]
\label{cc-signed-content}
For arbitrary integers \(x_k\), \(k\in K_B\), define the canonical product
\[
 Z_x=\prod_{x_k>0}w_k^{x_k}
           \prod_{x_k<0}\bar w_k^{-x_k},\qquad
 U_p(x)=\sum_k|x_k|v_p(Q_k),\qquad
 D_p(x)=\sum_k x_k d_{k,p}.
\]
The empty product is one. Then
\begin{align*}
 2v_p(\operatorname{cont}(Z_x))&=U_p(x)-|D_p(x)|,\\
 \log N((Z_x)_{\rm prim})&=\sum_p|D_p(x)|\log p.
\end{align*}
All sums have finite support.
\end{lemma}
\begin{proof}
Let \(E_p^+,E_p^-\) be the two ideal valuations of \(Z_x\). Their sum
is \(U_p(x)\), and their difference is \(D_p(x)\). A rational power
\(p^j\) divides both integer coordinates exactly when it divides both
conjugate ideal factors, so
\[
 v_p(\operatorname{cont}(Z_x))=\min(E_p^+,E_p^-).
\]
Now \(2\min(s,t)=s+t-|s-t|\) proves the first identity. Moreover
\[
 N(Z_x)=\prod_k Q_k^{|x_k|},\qquad
 N((Z_x)_{\rm prim})=
       N(Z_x)/\operatorname{cont}(Z_x)^2.
\]
Taking logarithms and using the first identity proves the second.
Neither the ramified prime three nor an inert prime occurs. The proof
includes an empty product and individual zero coordinates.
\end{proof}

\begin{lemma}[CC2: support of collective content]
\label{cc-content-cutoff}
For every subset \(J\subseteq K_B\), put \(W_J=\prod_{k\in J}w_k\).
Every prime dividing \(\operatorname{cont}(W_J)\) is at most \(12B-5\).
The same bound holds for the canonical signed products in
Lemma~\ref{cc-signed-content}.
\end{lemma}
\begin{proof}
A content prime occurs with both orientations. Primitivity of each
root forces occurrences at two distinct indices \(k,l\). If
\(p>12B-5\), then
\[
 p\mid Q_k-Q_l=3(k-l)\bigl(3(k+l)+1\bigr).
\]
But \(p>3\), \(0<|k-l|\le B-1<p\), and
\(0<3(k+l)+1\le12B-5<p\), a contradiction.
For a canonical signed product a prime above this cutoff can occur
at only one index and in only its selected orientation, giving the
same conclusion. Opposite occurrences of the same index must first
be cancelled: the statement is false for unrestricted products such
as \(w_k\bar w_k=Q_k\).
\end{proof}

\begin{theorem}[CC3: uniform interval content]
\label{cc-interval-content}
For every interval \(J=\{a,\ldots,a+H-1\}\subseteq K_B\), where
\(0\le H\le B\) is an integer, as \(B\to\infty\),
\begin{align*}
 \log\operatorname{cont}(W_J)
       &=\frac H2\log B+O(H\log\log B+B),\\
 \log N((W_J)_{\rm prim})
       &=H\log B+O(H\log\log B+B).
\end{align*}
The constants are absolute and independent of \(a,H\). In particular,
\begin{align*}
 \log\operatorname{cont}(W_{K_B})
       &=(1/2+o(1))B\log B,\\
 \log N((W_{K_B})_{\rm prim})
       &=(1+o(1))B\log B,
\end{align*}
whereas \(\log N(W_{K_B})=2B\log B+O(B)\).
\end{theorem}
\begin{proof}
The polynomial \(9X^2+3X+1\) has discriminant \(-27\). At each
split prime \(p\ge7\), its two roots modulo \(p\) are simple and each
lifts uniquely at every prime power. Label the two towers by the two
orientations. Since \(9B^2<Q_k<36B^2\), put
\[
 h_p=\left\lfloor\frac{\log(36B^2)}{\log p}\right\rfloor .
\]
No individual valuation exceeds \(h_p\). For either fixed tower, the
number of points in an interval of \(H\) consecutive integers belonging
to its residue class modulo \(p^j\) differs from \(H/p^j\) by at most
one. Therefore the total oriented valuation satisfies
\[
 \left|E_p^\pm(J)-\frac H{p-1}\right|\le h_p+1.
\]
Indeed the first \(h_p\) errors sum to at most \(h_p\), and the omitted
geometric tail is
\[
 \frac H{p^{h_p}(p-1)}<1,
\]
using \(p^{h_p+1}>36B^2\), \(H\le B\), and \(p\ge7\).
The minimum of the two valuations obeys the same error bound.
Lemma~\ref{cc-content-cutoff} then gives
\begin{align*}
 \log\operatorname{cont}(W_J)
 &=H\sum_{\substack{p\le12B\\p\equiv1\ (3)}}
       \frac{\log p}{p-1}
   +O\!\left(\sum_{p\le12B}(h_p+1)\log p\right).
\end{align*}
The endpoint sum is \(O(B)\): each summand is at most
\(\log(36B^2)+\log(12B)\), and the elementary Chebyshev bound gives
\(\pi(12B)=O(B/\log B)\). This estimate includes primes absent from
\(W_J\); their zero contributions still satisfy the oriented error bound.

The only progression-distribution input is the fixed-modulus estimate
\[
 \vartheta(x;3,1)=x/2+O(x/\log x).
\]
It follows from Theorem~1.2, equation~(1.12), printed page~5 of
\href{https://arxiv.org/abs/1802.00085v3}
{M.~A.~Bennett, G.~Martin, K.~O'Bryant and A.~Rechnitzer,
\emph{Explicit bounds for primes in arithmetic progressions}}.
That theorem gives explicit positive constants for each fixed modulus;
only modulus three is used here. Partial summation yields
\[
 \sum_{\substack{p\le x\\p\equiv1\ (3)}}\frac{\log p}{p}
       =\frac12\log x+O(\log\log x).
\]
The difference obtained by replacing \(1/p\) with \(1/(p-1)\) is
bounded by the convergent series \(\sum_{m\ge2}\log m/(m(m-1))\).
Substitution proves the first assertion. Finally
\(\log Q_k=2\log B+O(1)\) uniformly in \(K_B\), and
\[
 N(W_J)=\operatorname{cont}(W_J)^2 N((W_J)_{\rm prim})
\]
proves the remaining formulas. The asymptotic formulation avoids any
small-\(B\) convention for \(\log\log B\).
\end{proof}

\begin{lemma}[CC4: the exact deletion cost]
\label{cc-deletion}
If \(K_B=J\sqcup F\), then
\[
 \operatorname{cont}(W_{K_B})\mid
 \operatorname{cont}(W_J)\frac{N(W_F)}{\operatorname{cont}(W_F)}.
\]
The right side is an integer. Consequently
\[
 \log\operatorname{cont}(W_J)
 \ge\log\operatorname{cont}(W_{K_B})-\sum_{k\in F}\log Q_k.
\]
\end{lemma}
\begin{proof}
At a split prime write \((a,b)\) for the valuations of \(W_J\) and
\((c,d)\) for those of \(W_F\). Then
\[
 \min(a+c,b+d)\le\min(a,b)+\max(c,d),
 \qquad \max(c,d)=c+d-\min(c,d).
\]
These are exactly the primewise divisibility inequalities. They also
show that \(N(W_F)/\operatorname{cont}(W_F)\) is an integer. Taking
logarithms and discarding the nonnegative \(\log\operatorname{cont}(W_F)\)
proves the weaker displayed lower bound.
\end{proof}

Deleting \(o(B)\) arbitrary roots preserves the full-block main term.
Deleting a fixed positive fraction has a charged cost; this lemma does
not transfer that main term to the known independent domains of roughly
half the block. More fundamentally, the primes cancelled here divide
input norms \(Q_k\), whereas a boundary prime dividing
\(P(w_k^n)\) does not divide \(Q_k\). No formula above identifies the
cancelled input content with the negative boundary contributions
\((3-v_q(P(w_k^n)))_+\log q\). A prime hitting just one boundary need
not make the full-product ratio an \(n\)-torsion element, because the
unhit factors remain.

\subsection{Exact cancellation for a common multiplier and complement}

Write \(R(a,b)=(-b,a+b)\) and
\(\det((a,b),(c,d))=ad-bc\). Let \(W=(a,b)\) be primitive and let
\(V\ne0\) be an Eisenstein integer. Put
\[
 c=\operatorname{cont}(VW),\quad d=\operatorname{cont}(V),\quad
 U=VW/c,\quad V_0=V/d,
\]
and define
\[
 x=\det(U,V_0),\qquad D=\det(U,RV_0),\qquad
 y=\det(U,R^2V_0).
\]
Primitivity makes \(W\ne0\), so \(c,d\) are positive integers and
\(U,V_0\) are actual primitive pairs.

\begin{theorem}[CD1: the precise common scalar]
\label{cd-exact-scalar}
There is a positive integer \(r\) such that
\[
 N(V)=cdr,\qquad (x,D,y)=r(-b,a,a+b),
 \qquad \gcd(|x|,|D|,|y|)=r.
\]
\end{theorem}
\begin{proof}
Multiplication by \(V=v_0+v_1\zeta\) has integer matrix
\[
 \begin{pmatrix}v_0&-v_1\\v_1&v_0+v_1\end{pmatrix}
\]
of determinant \(N(V)\), and commutes with \(R\). Apply this determinant
identity to \(W\) and \(1,\zeta,\zeta^2\), and clear the positive
denominators \(c,d\). The result is
\[
 cdx=-bN(V),\qquad cdD=aN(V),\qquad cdy=(a+b)N(V).
\]
Choose integers \(u,v\) with \(ua+vb=1\). Then
\[
 N(V)=cd(uD-vx).
\]
Thus \(cd\mid N(V)\), and \(r=N(V)/(cd)\) is a positive integer.
Cancelling \(cd\) proves the three coordinates. Their common gcd is
\(r\), since \(\gcd(a,b,a+b)=1\). The proof permits individual zero
coordinates of \(W\); only its primitivity and \(V\ne0\) are needed.
In particular integrality of \(r\) follows from B\'ezout, not merely
from knowing that the three determinants are integers separately.
\end{proof}

\begin{corollary}[CD2: the primitive boundary ledger is unchanged]
\label{cd-ledger}
Assume \(ab(a+b)\ne0\). Dividing by the full common content in
Theorem~\ref{cd-exact-scalar} gives the additive relation
\(-b+(a+b)=a\). The absolute primitive boundary product is
\[
 |(x/r)(D/r)(y/r)|=|ab(a+b)|.
\]
Its radical and its complete signed depth sum are therefore exactly
those of \(W\):
\[
 \sum_{q\mid |(x/r)(D/r)(y/r)|}
       \bigl(v_q(|(x/r)(D/r)(y/r)|)-3\bigr)\log q
 =\sum_{q\mid |ab(a+b)|}\bigl(v_q(ab(a+b))-3\bigr)\log q.
\]
\end{corollary}
\begin{proof}
These identities follow from the three exact coordinates, including
their signs. Before division the determinant product is
\(-r^3ab(a+b)\). If \(r>1\), that triple is not primitive and cannot
be inserted into a primitive radical assertion before cancelling \(r\).
There is no omitted content term in the stated equalities.
\end{proof}

\begin{corollary}[CD3: simultaneous audit of every complement owner]
\label{cd-complements}
In the actual primitive unramified power block, let \(n\) be a positive
integer, and take any finite set \(J\) of indices and any \(k\in J\). Set
\[
 W=w_k^n,\qquad V=\prod_{l\in J\setminus\{k\}}w_l^n.
\]
The empty complement is one. The preceding identities hold for every
\(k\) simultaneously, and the primitive three-phase output is precisely
that owner's original additive triple.
\end{corollary}
\begin{proof}
The actual unramified primitive-power statement makes \(W\) primitive
with nonzero boundary arms, and every complement factor is nonzero.
Apply Theorem~\ref{cd-exact-scalar} and Corollary~\ref{cd-ledger}
separately to each \(k\); neither argument imposes a common boundary
prime condition across the other indices.
\end{proof}

This nonrectangular construction retains a different complement for
each owner but returns exactly the original primitive ledger. It audits
this particular common-multiplier map, without ruling out other mixed
determinants, divided differences, or additional relations between their
primitive outputs.

\subsection{The powerful part of the primitive collective input norm}

For an interval \(J\subseteq K_B\) of length \(H\), define
\[
 N_J=N((W_J)_{\rm prim}),\qquad
 R_J=\prod_{\substack{p>12B\\p\mid Q_k\text{ for some }k\in J}}p,
 \qquad
 M_B=\prod_{\substack{p\le12B\\p\equiv1\ (3)}}p^{h_p},
\]
where \(h_p=\lfloor\log(36B^2)/\log p\rfloor\). These are finite
products of actual rational primes.

\begin{theorem}[CP1: localization of the entire powerful part]
\label{cp-localization}
There is a positive integer \(D_J\) such that
\[
 N_J=R_JD_J,\qquad \gcd(R_J,D_J)=1,\qquad D_J\mid M_B.
\]
Moreover
\[
 \frac{N_J}{\operatorname{rad}(N_J)}\mid D_J\mid M_B,
 \qquad 0\le\log N_J-\log\operatorname{rad}(N_J)
                \le\log M_B=O(B).
\]
The implied constant is absolute; the divisibilities include an empty
interval.
\end{theorem}
\begin{proof}
For \(p\le12B\), let \(N_j^+,N_j^-\) count the two oriented residue
classes modulo \(p^j\) in \(J\). Each is either
\(\lfloor H/p^j\rfloor\) or \(\lceil H/p^j\rceil\). Thus
\( |N_j^+-N_j^-|\le1\), and Lemma~\ref{cc-signed-content} gives
\[
 v_p(N_J)=|E_p^+-E_p^-|
    =\left|\sum_{j=1}^{h_p}(N_j^+-N_j^-)\right|\le h_p.
\]
If \(p>12B\), the difference-of-norms argument in
Lemma~\ref{cc-content-cutoff} shows that it divides at most one norm
in the entire block. Its depth is one because
\(p^2>(12B)^2>36B^2>Q_k\). Its sole orientation survives primitive
cancellation. These are exactly the primes of \(R_J\). No inert or
ramified primes occur. The actual prime factorization proves all the
divisibilities. Finally the Chebyshev bound gives
\[
 \log M_B\le\pi(12B)\log(36B^2)=O(B).
\]
No prime number theorem in a progression is used for this last bound.
\end{proof}

\begin{corollary}[CP2: radical mass of the primitive input norm]
\label{cp-radical-mass}
For each fixed \(\delta>0\), uniformly on intervals with \(H\ge\delta B\),
as \(B\to\infty\),
\begin{align*}
 \log N_J&=H\log B+O(H\log\log B+B),\\
 \log R_J&=H\log B+O(H\log\log B+B),\\
 \frac{\log\operatorname{rad}(N_J)}{\log N_J}
          &\ge1-O_\delta(1/\log B).
\end{align*}
\end{corollary}
\begin{proof}
The first formula is Theorem~\ref{cc-interval-content}. Subtracting
\(\log D_J=O(B)\) proves the second. For all sufficiently large \(B\)
depending only on \(\delta\),
\(\log N_J\ge(\delta/2)B\log B\). Divide the logarithmic powerful-part
bound of Theorem~\ref{cp-localization} by this lower bound to obtain
the last assertion. The constants in the first two errors are absolute.
\end{proof}

\begin{theorem}[CP3: obstruction to compressing this aggregate]
\label{cp-compression}
Suppose positive integers \(V,Q,g\) satisfy
\[
 N_J=VQ^g,\qquad Q>1,\qquad g\ge2.
\]
Then
\[
 R_J\mid V,\qquad Q^g\mid D_J,\qquad
 \frac{\log V}{g}\ge
     \frac{(\log2)(\log R_J)}{\log D_J}.
\]
Here \(D_J>1\). For every fixed \(\delta>0\), there are positive
\(c_\delta,B_\delta\) such that for every integer \(B\ge B_\delta\),
every interval \(H\ge\delta B\), and every such representation,
\[
 \frac{\max(1,\log V)}g\ge c_\delta\log B.
\]
No power-free condition on \(V\), or coprimality of \(V,Q\), is assumed.
\end{theorem}
\begin{proof}
Every prime of \(R_J\) has total depth one. Since \(g\ge2\), its
exponent in \(Q\) must be zero and its exponent in \(V\) must be one.
All prime factors of \(Q\) instead belong to \(D_J\), and their
\(g\)-fold exponents are bounded by the corresponding exponents in
\(D_J\). This proves the two divisibilities. Since \(Q\ge2\),
\[
 g\log2\le g\log Q\le\log D_J.
\]
Combine this with \(\log V\ge\log R_J\). For the uniform conclusion,
use \(\log D_J\le CB\) with one absolute \(C\), and
\(\log R_J\ge(\delta/2)B\log B\) for sufficiently large \(B\).
If \(D_J=1\), no representation with \(Q>1\) exists.
\end{proof}

\begin{theorem}[CP4: a primewise exponent ceiling and linear obstruction]
\label{cp-linear-obstruction}
Under the hypotheses of Theorem~\ref{cp-compression},
\[
 g\le\frac{\log(36B^2)}{\log7},\qquad
 \frac{\log V}{g}\ge
       \frac{(\log7)(\log R_J)}{\log(36B^2)}.
\]
Consequently, for every fixed \(\delta>0\), all sufficiently large
integers \(B\), all intervals \(H\ge\delta B\), and every representation
in that theorem,
\[
 \frac{\max(1,\log V)}g\ge\frac{\delta\log7}{6}B.
\]
\end{theorem}
\begin{proof}
Choose an actual prime \(p\mid Q\), possible because \(Q>1\).
It divides \(D_J\), is split, and is at least seven. The actual exponent
comparison and Theorem~\ref{cp-localization} give
\[
 g\le g v_p(Q)\le v_p(D_J)\le h_p
             \le\frac{\log(36B^2)}{\log7}.
\]
Together with \(\log V\ge\log R_J\), this proves both finite bounds.
When \(\log R_J=0\), the second bound is still the valid zero lower
bound. On the specified macroscopic intervals,
\(\log R_J\ge(\delta/2)B\log B\) for sufficiently large \(B\), and
\(B\ge36\) gives \(\log(36B^2)\le3\log B\). Substitution proves the
linear lower bound. The ceiling uses a single actual prime of \(Q\),
not merely the total mass \(\log D_J\).
\end{proof}

The height decrease of the primitive collective product is genuine, and
almost all of its input norm mass is radical mass. Nevertheless this
specific object cannot produce small-\(\lambda\) compressed norm seeds
on macroscopic intervals: every possible displayed representation has
\(\lambda\) growing at least linearly in \(B\). This excludes neither
useful representations of original roots nor a different combined map.

The proofs in this section are ordinary proofs. The sole progression
input has been stated with its primary theorem and exact fixed-modulus
scope. Separate finite replays of the content identities, complement
identities, and norm profiles test exact instances; they supply no
evidence in place of the asymptotic proofs or their universal quantifiers.
This section makes no additional Lean claim. Any separately formalized
arithmetic core has its own explicit signatures and verification scope.
In particular the unrestricted boundary signed tail, individual exceptional
roots, the independent-domain complement, and general ABC coverage remain
open.
